% Test report -- example for the plaindoc class. All data are fictional.
% Build: latexmk -lualatex dp305-thermal-report.tex
\documentclass[type=report, lang=en-GB]{plaindoc}
\usepackage[style=numeric-comp, sorting=none, maxbibnames=3]{biblatex}
\addbibresource{refs.bib}

\pdset{
  org      = {Orrin Instruments · Product Engineering (fictional)},
  title    = {DP305 thermal derating at full load},
  subtitle = {Test report},
  author   = {M. Chan, R. Okafor},
  date     = {25 November 2026},
  docid    = {TR-2026-041},
  version  = {Rev. B},
  status   = {Released},
}

\begin{document}
\maketitle

\begin{summary}
At \qty{25}{\degreeCelsius} ambient the DP305 delivered \qty{30}{\volt} at
\qty{5}{\ampere} for \qty{4}{\hour} without tripping over-temperature
protection (OTP); the heatsink settled at \qty{71.4}{\degreeCelsius}. At
\qty{40}{\degreeCelsius} ambient, OTP switched the output off after
\qty{65}{\minute}. We recommend specifying full output up to
\qty{35}{\degreeCelsius} and derating the current to \qty{4}{\ampere} at
\qty{40}{\degreeCelsius}.
\end{summary}

\section{Scope}

The DP305 user manual~\cite{dp305um} rates the supply for operation from
\qtyrange{0}{40}{\degreeCelsius} but does not state whether full output is
available across that range. This test measured the heatsink temperature at
full output in three ambient temperatures and recorded whether OTP, which
trips at \qty{85}{\degreeCelsius}, was reached.

The test does not cover component lifetime or the touch-temperature limits
of \mbox{IEC 61010-1}~\cite{iec61010}; those are assessed in TR-2026-044.

\section{Method}

One production unit (serial number 3051187) was placed in a climatic chamber
and run at constant ambient temperature following the dry-heat procedure of
\mbox{IEC 60068-2-2}~\cite{iec60068}. The output was set to \qty{30.00}{\volt} with a
\qty{5.000}{\ampere} limit and loaded with an electronic load in constant-current
mode. \Cref{tab:equipment} lists the equipment.

\begin{table}
\caption{Test equipment.}
\label{tab:equipment}
\begin{tabularx}{\linewidth}{@{}lL l@{}}
\toprule
\hd{Item} & \hd{Model and use} & \hd{Calibrated} \\
\midrule
Climatic chamber   & CT-150, ambient set point ±\qty{0.5}{\kelvin}       & 2026-03 \\
Electronic load    & EL-600, constant current \qty{5.000}{\ampere}        & 2026-05 \\
Thermocouples      & Type K, \num{4} channels, bonded with thermal epoxy  & 2026-05 \\
Data logger        & DL-16, \qty{10}{\s} sample interval                  & 2026-05 \\
\bottomrule
\end{tabularx}
\end{table}

Thermocouples were placed on the heatsink next to the output transistor
(TC1), on the transformer winding (TC2), at the air outlet (TC3) and in the
chamber (TC4). Each run lasted \qty{4}{\hour} or until OTP tripped. The
chamber was held at the set point for \qty{1}{\hour} before each run.

\section{Results}

\Cref{fig:heatsink} shows the heatsink temperature (TC1) for each ambient
temperature. At \qty{25}{\degreeCelsius} and \qty{35}{\degreeCelsius} the
temperature reached a steady state within \qty{90}{\minute}. At
\qty{40}{\degreeCelsius} it rose past the OTP threshold, the output switched
off at \qty{65}{\minute}, and the heatsink cooled towards ambient.

\begin{figure}
\centering
\begin{tikzpicture}
\begin{axis}[pd/axis, width=.8\linewidth, height=62mm, clip=false,
    xmin=0, xmax=240, ymin=20, ymax=95, xtick={0,30,...,240},
    xlabel={Time (min)}, ylabel={Heatsink temperature (°C)}]
  \addplot+[domain=0:240, samples=49, mark repeat=4] {25+46.4*(1-exp(-x/22))};
  \addplot+[domain=0:240, samples=49, mark repeat=4] {35+47.0*(1-exp(-x/22))};
  \addplot+[domain=0:65, samples=14, mark repeat=4] {40+47.5*(1-exp(-x/22))};
  \addplot+[oi-green, dotted, no marks, domain=65:240, samples=36, forget plot]
    {40+45.0*exp(-(x-65)/30)};
  \draw[pd-ink2, line width=.4pt, dash pattern=on 3pt off 2pt] (axis cs:0,85) -- (axis cs:240,85);
  \node[anchor=south, font=\sffamily\scriptsize, text=pd-ink2] at (axis cs:165,85.5) {OTP threshold \qty{85}{\degreeCelsius}};
  \node[anchor=south west, font=\sffamily\scriptsize] at (axis cs:67,86) {OTP trips};
  \node[anchor=west, font=\sffamily\scriptsize] at (axis cs:244,71.4) {\qty{25}{\degreeCelsius} ambient};
  \node[anchor=west, font=\sffamily\scriptsize] at (axis cs:244,82.0) {\qty{35}{\degreeCelsius} ambient};
  \node[anchor=west, font=\sffamily\scriptsize, align=left] at (axis cs:244,41.5) {\qty{40}{\degreeCelsius} ambient\\[-1pt](output off)};
\end{axis}
\end{tikzpicture}
\caption{Heatsink temperature (TC1) at \qty{30}{\volt}, \qty{5}{\ampere}.
  After OTP trips the output is off and the heatsink cools.}
\label{fig:heatsink}
\end{figure}

\Cref{tab:steady} gives the temperatures at the end of each run, or at the
moment OTP tripped.

\begin{table}
\caption{Temperatures at the end of each run (°C).}
\label{tab:steady}
\begin{tabular*}{\linewidth}{@{\extracolsep{\fill}}lS[table-format=2.1]S[table-format=2.1]S[table-format=2.1]l@{}}
\toprule
\hd{Ambient} & {\hd{Heatsink (TC1)}} & {\hd{Transformer (TC2)}} & {\hd{Outlet air (TC3)}} & \hd{OTP} \\
\midrule
\qty{25}{\degreeCelsius} & 71.4 & 64.8 & 41.2 & No \\
\qty{35}{\degreeCelsius} & 82.0 & 75.1 & 50.9 & No \\
\qty{40}{\degreeCelsius} & 85.0 & 79.6 & 55.3 & After \qty{65}{\minute} \\
\bottomrule
\end{tabular*}
\end{table}

\section{Discussion}

The heatsink sits about \qty{47}{\kelvin} above ambient at full load, in
every run. With OTP at \qty{85}{\degreeCelsius}, full output is therefore
sustainable up to about \qty{38}{\degreeCelsius} ambient; the
\qty{35}{\degreeCelsius} run leaves a margin of \qty{3}{\kelvin}.
Dissipation in the output stage scales with output current, so reducing
the current to \qty{4}{\ampere} lowers the temperature rise to roughly
\qty{38}{\kelvin}\footnote{Estimated from the measured rise of
\qty{47}{\kelvin} at \qty{5}{\ampere}, assuming linear dependence on current
at constant output voltage. A confirmation run is planned.} and keeps the
heatsink below the threshold at \qty{40}{\degreeCelsius}.

One unit was tested. Unit-to-unit variation in fan speed and heatsink
contact could move the threshold by a few kelvin either way.

\section{Recommendations}

\begin{enumerate}
\item State in the specifications that full output is available up to
      \qty{35}{\degreeCelsius} ambient.
\item Derate the output current linearly from \qty{5}{\ampere} at
      \qty{35}{\degreeCelsius} to \qty{4}{\ampere} at \qty{40}{\degreeCelsius}.
\item Repeat the \qty{40}{\degreeCelsius} run at \qty{4}{\ampere} on three
      units before the datasheet is released.
\end{enumerate}

\printbibliography[heading=bibnumbered, title=References]

\end{document}
